% The title and chapter boundary remain provisional. This chapter is the
% complete working manuscript for the second seminar talk. Sections marked as
% supplements are intended for the written notes, not for the default live
% route through the talk.

\chapter{Smooth geometry through \texorpdfstring{$\Cinfty$}{C-infinity}-rings}
\label[chapter]{chap:Smooth_Geometry_Through_Cinfty_Rings}

\section{Smooth functional calculus}
\label[section]{sec:Smooth_Functional_Calculus}

The preceding chapter ended with a basic loss of information. The equations
\[
  t=0
  \qquad\text{and}\qquad
  t^2=0
\]
have the same set of solutions in $\R$, but their linearizations at the origin
are different. Passing immediately to the zero set forgets the equation which
cut it out.

Ordinary algebraic geometry suggests retaining a quotient of a function ring.
In the smooth setting, the two candidates are
\[
  \Cinfty(\R)/(t)
  \qquad\text{and}\qquad
  \Cinfty(\R)/(t^2).
\]
These expressions lead to two questions.

\begin{enumerate}
  \item What structure on $\Cinfty(\R)$ remembers substitution by arbitrary
    smooth functions, rather than only addition and multiplication?
  \item Why does that structure descend to a quotient by an ordinary ideal?
\end{enumerate}

The answer to the first question is the notion of a $\Cinfty$-ring. The answer
to the second is Hadamard's lemma. Once both are in place, we will compute the
two displayed quotients and see exactly what the second remembers beyond its
single real point.

There is also a more global test for the language. If a manifold is replaced
by its ring of smooth functions, can we still recover its smooth maps? The
principal geometric result of the chapter says that we can. In fact, the same
argument works for arbitrary closed subsets of Euclidean spaces equipped with
their smooth functions.

Let us begin with the motivating example. If $M$ is a manifold, a smooth map
$f\colon\R^n\to\R$ can be applied pointwise to an $n$-tuple of smooth
functions on $M$:
\[
  (a_1,\ldots,a_n)
  \longmapsto
  f(a_1,\ldots,a_n),
  \qquad
  f(a_1,\ldots,a_n)(x)
  =f(a_1(x),\ldots,a_n(x)).
\]
A $\Cinfty$-ring is a set on which all these smooth operations make sense,
subject to the relations imposed by composition of smooth maps.
The operations formulation below follows Moerdijk--Reyes
\cite[Chapter~I, Section~1]{Moerdijk_Reyes_Smooth_Infinitesimal_Analysis}; a
modern account, including the functor formulation, is given by Joyce
\cite[Sections~2.1--2.2]{Joyce_Introduction_Cinfty_Schemes}.

\begin{definition}[$\Cinfty$-ring]
  \label[definition]{def:Cinfty_Ring_Operations}
  A \emph{$\Cinfty$-ring} is a set $A$ equipped, for every $n\geq 0$ and every
  smooth map $f\colon\R^n\to\R$, with an operation
  \[
    \Phi_f\colon A^n\longrightarrow A.
  \]
  These operations satisfy the following two conditions.

  \begin{enumerate}
    \item If $\pi_i\colon\R^n\to\R$ is the $i$th projection, then
      $\Phi_{\pi_i}(a_1,\ldots,a_n)=a_i$.
    \item If $f_1,\ldots,f_m\colon\R^n\to\R$ and
      $g\colon\R^m\to\R$ are smooth, then
      \[
        \Phi_{g\circ(f_1,\ldots,f_m)}
        =\Phi_g\circ(\Phi_{f_1},\ldots,\Phi_{f_m}).
      \]
  \end{enumerate}
\end{definition}

The case $n=0$ supplies an element of $A$ for every constant in $\R$. The
second axiom says that evaluating a composite smooth function gives the same
answer as evaluating its pieces successively.

\begin{example}[Smooth functions]
  \label[example]{ex:Smooth_Functions_Cinfty_Ring}
  For every manifold $M$, the set $\Cinfty(M)$ is a $\Cinfty$-ring under
  pointwise smooth functional calculus. More generally, the same construction
  works whenever a notion of smooth real-valued function has been specified
  and is closed under smooth substitution.
\end{example}

\begin{definition}[$\Cinfty$-homomorphism]
  \label[definition]{def:Cinfty_Homomorphism}
  A homomorphism of $\Cinfty$-rings is a map $\varphi\colon A\to B$ satisfying
  \[
    \varphi\bigl(\Phi_f(a_1,\ldots,a_n)\bigr)
    =\Phi_f\bigl(\varphi(a_1),\ldots,\varphi(a_n)\bigr)
  \]
  for every smooth $f\colon\R^n\to\R$ and every
  $(a_1,\ldots,a_n)\in A^n$.
\end{definition}

If $u\colon M\to N$ is smooth, then pullback gives a homomorphism
\[
  u^*\colon\Cinfty(N)\longrightarrow\Cinfty(M),
  \qquad
  h\longmapsto h\circ u.
\]
Thus smooth geometry enters the algebra contravariantly. This reversal of
arrows will remain in force throughout the next several chapters.

\begin{observation}[The underlying commutative algebra]
  \label[observation]{obs:Underlying_Commutative_Algebra}
  Every $\Cinfty$-ring has an underlying commutative $\R$-algebra. Addition,
  multiplication, scalar multiplication, zero, and one are obtained from the
  corresponding smooth maps between Euclidean spaces. Every
  $\Cinfty$-homomorphism is consequently a homomorphism of commutative
  $\R$-algebras.
\end{observation}

The converse is false. The extra smooth operations impose genuine algebraic
conditions. For example, in every $\Cinfty$-ring and for every $a\in A$, the
element $1+a^2$ is invertible. Its inverse is
\[
  \Phi_g(a),
  \qquad
  g(x)=\frac{1}{1+x^2},
\]
because the identity $(1+x^2)g(x)=1$ is an identity of smooth functions.
Therefore the usual commutative algebra $\R[t]$ cannot be the underlying
algebra of a $\Cinfty$-ring whose distinguished element $t$ is its polynomial
generator: the polynomial $1+t^2$ has no inverse in $\R[t]$.

There is a compact categorical formulation of the same definition. Let
$\mathrm{Euc}$ be the category whose objects are the Euclidean spaces
$\R^n$, for $n\geq 0$, and whose morphisms are smooth maps. Its finite
products are the usual products of Euclidean spaces.

\begin{proposition}[Functor formulation]
  \label[proposition]{prop:Cinfty_Ring_Functor_Formulation}
  To give a $\Cinfty$-ring is equivalently to give a finite-product-preserving
  functor
  \[
    A\colon\mathrm{Euc}\longrightarrow\mathrm{Set}.
  \]
  Under this correspondence, the underlying set of the ring is $A(\R)$.
\end{proposition}

\begin{proof}
  Suppose first that the operations of
  \Cref{def:Cinfty_Ring_Operations} are given. Set
  $A(\R^n)=A^n$. If
  $f=(f_1,\ldots,f_m)\colon\R^n\to\R^m$, define
  \[
    A(f)=(\Phi_{f_1},\ldots,\Phi_{f_m})\colon A^n\longrightarrow A^m.
  \]
  The projection and composition axioms say precisely that this is a functor,
  and the construction visibly preserves finite products.

  Conversely, let $A\colon\mathrm{Euc}\to\mathrm{Set}$ preserve finite
  products and put $A=A(\R)$. The product comparison identifies
  $A(\R^n)$ with $A^n$. For a smooth map $f\colon\R^n\to\R$, the map $A(f)$
  then defines $\Phi_f\colon A^n\to A$. Functoriality gives the two axioms in
  \Cref{def:Cinfty_Ring_Operations}.
\end{proof}

This is an instance of a Lawvere theory, but no general theory of Lawvere
theories is needed here. The useful point is simply that all smooth maps
between Euclidean spaces act functorially.

\section{Free smooth algebras and singular equations}
\label[section]{sec:Free_Smooth_Algebras_And_Singular_Equations}

\subsection{Free smooth algebras}
\label[subsection]{sec:Free_Smooth_Algebras}

Let $x_1,\ldots,x_n\in\Cinfty(\R^n)$ be the coordinate functions. Every
smooth function $f\in\Cinfty(\R^n)$ is obtained by applying the operation
$\Phi_f$ to these coordinates:
\[
  f=\Phi_f(x_1,\ldots,x_n).
\]
This tautological identity is the key to the universal property.

\begin{theorem}[Free $\Cinfty$-rings]
  \label[theorem]{thm:Free_Cinfty_Rings}
  For every $\Cinfty$-ring $A$, evaluation on the coordinate functions induces
  a bijection
  \[
    \Hom_{\Cinfty\text{-}\mathrm{Ring}}
    \bigl(\Cinfty(\R^n),A\bigr)
    \longrightarrow A^n,
    \qquad
    \varphi\longmapsto
    \bigl(\varphi(x_1),\ldots,\varphi(x_n)\bigr).
  \]
  Thus $\Cinfty(\R^n)$ is the free $\Cinfty$-ring on
  $x_1,\ldots,x_n$.
\end{theorem}

This is Moerdijk--Reyes, Proposition~1.1
\cite[Chapter~I, Proposition~1.1]{Moerdijk_Reyes_Smooth_Infinitesimal_Analysis}.

\begin{proof}
  Given $(a_1,\ldots,a_n)\in A^n$, define
  \[
    \operatorname{ev}_{a_1,\ldots,a_n}
    \colon\Cinfty(\R^n)\longrightarrow A,
    \qquad
    f\longmapsto\Phi_f(a_1,\ldots,a_n).
  \]
  The composition axiom for smooth functional calculus says exactly that this
  map preserves every smooth operation, so it is a $\Cinfty$-homomorphism. It
  sends $x_i$ to $a_i$.

  Conversely, if $\varphi\colon\Cinfty(\R^n)\to A$ is a
  $\Cinfty$-homomorphism, then
  \[
    \varphi(f)
    =\varphi\bigl(\Phi_f(x_1,\ldots,x_n)\bigr)
    =\Phi_f\bigl(\varphi(x_1),\ldots,\varphi(x_n)\bigr).
  \]
  Hence $\varphi$ is uniquely determined by the images of the coordinate
  functions, and the two constructions are inverse.
\end{proof}

Smooth functions on $\R^n$ therefore play the role that polynomials in $n$
variables play in ordinary commutative algebra. The analogy is useful, but it
should not be taken to say that a smooth function is determined by a finite
algebraic expression. The allowed operations already include every smooth map
$\R^n\to\R$.

\subsection{Quotients and smooth equations}
\label[subsection]{sec:Quotients_And_Smooth_Equations}

To use quotient rings, we must show that an ordinary ideal is respected by all
smooth operations. The required fact is a multivariable form of the
fundamental theorem of calculus.

\begin{lemma}[Hadamard's lemma]
  \label[lemma]{lem:Hadamard}
  For every smooth map $f\colon\R^n\to\R$, there are smooth maps
  $g_1,\ldots,g_n\colon\R^n\times\R^n\to\R$ such that
  \[
    f(y)-f(x)=\sum_{i=1}^n(y_i-x_i)g_i(x,y).
  \]
\end{lemma}

\begin{proof}
  Put
  \[
    g_i(x,y)
    =\int_0^1
      \frac{\partial f}{\partial x_i}\bigl(x+s(y-x)\bigr)\,ds.
  \]
  Differentiation under the integral sign shows that $g_i$ is smooth. The
  fundamental theorem of calculus applied to
  $s\mapsto f(x+s(y-x))$ gives
  \[
    f(y)-f(x)
    =\int_0^1\frac{d}{ds}f(x+s(y-x))\,ds
    =\sum_{i=1}^n(y_i-x_i)g_i(x,y).
  \]
\end{proof}

\begin{proposition}[Quotient $\Cinfty$-rings]
  \label[proposition]{prop:Quotient_Cinfty_Rings}
  Let $A$ be a $\Cinfty$-ring, and let $I\subseteq A$ be an ideal in its
  underlying commutative algebra. There is a unique $\Cinfty$-ring structure
  on $A/I$ for which the quotient map $A\to A/I$ is a
  $\Cinfty$-homomorphism.
\end{proposition}

This is Moerdijk--Reyes, Proposition~1.2
\cite[Chapter~I, Proposition~1.2]{Moerdijk_Reyes_Smooth_Infinitesimal_Analysis}.

\begin{proof}
  The only possible definition is
  \[
    \Phi_f(\overline{a}_1,\ldots,\overline{a}_n)
    =\overline{\Phi_f(a_1,\ldots,a_n)}.
  \]
  We check that it is independent of the representatives. Suppose that
  $b_i-a_i\in I$ for every $i$. Apply \Cref{lem:Hadamard} inside $A$ to obtain
  \[
    \Phi_f(b_1,\ldots,b_n)-\Phi_f(a_1,\ldots,a_n)
    =\sum_{i=1}^n(b_i-a_i)
      \Phi_{g_i}(a_1,\ldots,a_n,b_1,\ldots,b_n).
  \]
  The right-hand side lies in $I$. Thus the operations are well defined. Their
  projection and composition axioms descend from those in $A$. Uniqueness is
  forced by the requirement that the quotient map preserve every operation.
\end{proof}

This result is special to smooth functional calculus. It says that every
ordinary ideal in the underlying algebra is already compatible with all the
smooth operations.

\begin{definition}[Finite generation and finite presentation]
  \label[definition]{def:Cinfty_Finiteness}
  A $\Cinfty$-ring $A$ is \emph{finitely generated} if there are
  $a_1,\ldots,a_n\in A$ such that no proper $\Cinfty$-subring contains all the
  $a_i$. It is \emph{finitely presented} if it is isomorphic to
  \[
    \Cinfty(\R^n)/(f_1,\ldots,f_k)
  \]
  for some finite collections of generators and relations.
\end{definition}

\begin{proposition}
  \label[proposition]{prop:Finitely_Generated_Cinfty_Quotient}
  A $\Cinfty$-ring is finitely generated if and only if it is isomorphic to
  $\Cinfty(\R^n)/I$ for some $n$ and some ideal
  $I\subseteq\Cinfty(\R^n)$.
\end{proposition}

\begin{proof}
  If $a_1,\ldots,a_n$ generate $A$, the universal property in
  \Cref{thm:Free_Cinfty_Rings} gives a homomorphism
  $\Cinfty(\R^n)\to A$ sending $x_i$ to $a_i$. Its image is a
  $\Cinfty$-subring containing the generators, so it is surjective. By
  \Cref{prop:Quotient_Cinfty_Rings}, the induced isomorphism of commutative
  algebras $\Cinfty(\R^n)/I\to A$ is an isomorphism of $\Cinfty$-rings. The
  converse follows because the images of the coordinate functions generate
  every such quotient.
\end{proof}

Finite generation here means generation under every smooth operation. It does
not imply that the underlying commutative $\R$-algebra is finitely generated.
Likewise, finite presentation is a genuine additional condition: smooth
function rings are not Noetherian. The next chapter will return to finite
presentation for manifolds.

\subsection{One point, two equations}
\label[subsection]{sec:One_Point_Two_Equations}

We can now compute the two quotients from the opening. In one variable,
\Cref{lem:Hadamard} gives
\[
  f(t)=f(0)+t g(t)
\]
for some smooth $g$. Evaluation at the origin therefore induces an
isomorphism
\[
  \Cinfty(\R)/(t)\cong\R.
\]

For the second quotient, Taylor's formula with smooth remainder gives
\[
  f(t)=f(0)+f'(0)t+t^2h(t).
\]
Consequently, the map
\[
  \Cinfty(\R)\longrightarrow\R[\varepsilon]/(\varepsilon^2),
  \qquad
  f\longmapsto f(0)+f'(0)\varepsilon
\]
is a surjective homomorphism of commutative $\R$-algebras with kernel $(t^2)$.
It identifies
\[
  \Cinfty(\R)/(t^2)
  \cong
  \R[\varepsilon]/(\varepsilon^2).
\]
The ring on the right is called the ring of \emph{dual numbers}. Its
$\Cinfty$-structure is not additional arbitrary data. It is the quotient
structure from \Cref{prop:Quotient_Cinfty_Rings}, and its operations are given
by the first-order Taylor formula
\[
  f(a+b\varepsilon)=f(a)+b f'(a)\varepsilon.
\]

To compare these quotients with their ordinary solution sets, we isolate the
appropriate notion of point.

\begin{definition}[Real point]
  \label[definition]{def:Real_Point}
  A \emph{real point} of a $\Cinfty$-ring $A$ is a
  $\Cinfty$-homomorphism $A\to\R$, where $\R$ carries its evident smooth
  functional calculus.
\end{definition}

\begin{proposition}[Real points of a quotient]
  \label[proposition]{prop:Real_Points_Of_Quotient}
  For every ideal $I\subseteq\Cinfty(\R^n)$, evaluation induces a bijection
  \[
    \Hom_{\Cinfty\text{-}\mathrm{Ring}}
    \bigl(\Cinfty(\R^n)/I,\R\bigr)
    \cong
    \{p\in\R^n\mid f(p)=0\text{ for every }f\in I\}.
  \]
\end{proposition}

\begin{proof}
  By \Cref{thm:Free_Cinfty_Rings}, every homomorphism
  $\Cinfty(\R^n)\to\R$ is evaluation at a unique point
  $p=(p_1,\ldots,p_n)$. It factors through the quotient precisely when every
  element of $I$ vanishes at $p$.
\end{proof}

Both $\Cinfty(\R)/(t)$ and $\Cinfty(\R)/(t^2)$ therefore have exactly one real
point. They are nevertheless not isomorphic: the class of $t$ is zero in the
first and is a nonzero square-zero element in the second. The quotient by
$(t^2)$ remembers first-order information which the zero set forgets.

\begin{warning}
  \label[warning]{warn:Ordinary_Cinfty_Not_Derived}
  This example does not yet construct a derived zero locus. An ordinary
  quotient $\Cinfty$-ring can retain nilpotents and dependence on the chosen
  ideal, but it does not by itself encode homotopy-coherent equations or
  derived pullbacks. It is the classical algebraic input from which the later
  derived theory will be built.
\end{warning}

\section{Recovering smooth spaces from functions}
\label[section]{sec:Recovering_Smooth_Spaces_From_Functions}

\subsection{Recovering closed smooth subsets}
\label[subsection]{sec:Recovering_Closed_Smooth_Subsets}

We now turn from equations to geometry. Let $X\subseteq\R^n$ be an arbitrary
closed subset. It need not be a manifold.

\begin{definition}[Smooth functions and maps on subsets]
  \label[definition]{def:Smooth_Functions_On_Subsets}
  A function $a\colon X\to\R$ is \emph{smooth} if it extends to a smooth
  function on an open neighborhood of $X$ in $\R^n$. We write $\Cinfty(X)$
  for the resulting $\Cinfty$-ring.

  A map $u\colon X\to Y\subseteq\R^m$ is \emph{smooth} if its component
  functions are smooth on $X$. Equivalently, it extends to a smooth map into
  $\R^m$ on an open neighborhood of $X$.
\end{definition}

The equivalence in the second paragraph uses only that there are finitely many
component functions: intersect their extension neighborhoods and combine the
extensions.

\begin{lemma}[Smooth extension from a closed subset]
  \label[lemma]{lem:Smooth_Extension_Closed_Subset}
  If $X\subseteq\R^n$ is closed, every element of $\Cinfty(X)$ is the
  restriction of a globally defined smooth function on $\R^n$.
\end{lemma}

\begin{proof}
  Let $a\in\Cinfty(X)$, and choose a smooth extension $\widetilde a$ on an open
  neighborhood $U$ of $X$. A smooth partition of unity subordinate to the
  cover
  \[
    \{U,\R^n\setminus X\}
  \]
  supplies a smooth function $\chi\colon\R^n\to[0,1]$ whose support is
  contained in $U$ and whose restriction to $X$ is one. Define
  \[
    \overline a(x)=
    \begin{cases}
      \chi(x)\widetilde a(x),&x\in U,\\
      0,&x\notin U.
    \end{cases}
  \]
  Since the support of $\chi$ is contained in $U$, this formula defines a
  smooth function on all of $\R^n$. Its restriction to $X$ is $a$.
\end{proof}

Let
\[
  \mathfrak m_X
  =\{h\in\Cinfty(\R^n)\mid h|_X=0\}.
\]
The restriction homomorphism is surjective by
\Cref{lem:Smooth_Extension_Closed_Subset}, with kernel $\mathfrak m_X$.
Hence
\[
  \Cinfty(X)\cong\Cinfty(\R^n)/\mathfrak m_X.
\]
In particular, $\Cinfty(X)$ is a finitely generated $\Cinfty$-ring, even when
$X$ is singular or has complicated topology.

The following theorem is the main geometric payoff of the chapter.

\begin{theorem}[Full faithfulness for closed subsets]
  \label[theorem]{thm:Full_Faithfulness_Closed_Subsets}
  Let $X\subseteq\R^n$ and $Y\subseteq\R^m$ be closed subsets. Pullback induces
  a bijection
  \[
    \{\text{smooth maps }Y\to X\}
    \longrightarrow
    \Hom_{\Cinfty\text{-}\mathrm{Ring}}
    \bigl(\Cinfty(X),\Cinfty(Y)\bigr).
  \]
  Consequently, the functor
  \[
    \Cinfty(-)\colon
    (\text{closed subsets of Euclidean spaces})\catop
    \longrightarrow
    \Cinfty\text{-}\mathrm{Ring}
  \]
  is fully faithful.
\end{theorem}

\begin{proof}
  Let
  \[
    \Psi\colon\Cinfty(X)\longrightarrow\Cinfty(Y)
  \]
  be a homomorphism. Write $x_1,\ldots,x_n$ for the coordinate functions on
  $\R^n$, and define
  \[
    g_i=\Psi(x_i|_X)\in\Cinfty(Y),
    \qquad
    g=(g_1,\ldots,g_n)\colon Y\longrightarrow\R^n.
  \]
  The map $g$ is smooth by \Cref{def:Smooth_Functions_On_Subsets}. For every
  $h\in\Cinfty(\R^n)$, preservation of smooth functional calculus gives
  \[
    \Psi(h|_X)
    =\Psi\bigl(\Phi_h(x_1|_X,\ldots,x_n|_X)\bigr)
    =\Phi_h(g_1,\ldots,g_n)
    =h\circ g.
  \]

  We claim that $g$ lands in $X$. If $g(y)=p\notin X$ for some $y\in Y$,
  choose a smooth bump function $h\colon\R^n\to\R$ supported in
  $\R^n\setminus X$ and satisfying $h(p)=1$. Then $h|_X=0$, whereas the
  displayed identity gives
  \[
    0=\Psi(h|_X)(y)=h(g(y))=1,
  \]
  a contradiction.

  It remains to identify $\Psi$. Given $a\in\Cinfty(X)$, choose a global
  smooth extension $h\in\Cinfty(\R^n)$ using
  \Cref{lem:Smooth_Extension_Closed_Subset}. Then
  \[
    \Psi(a)=\Psi(h|_X)=h\circ g=a\circ g.
  \]
  Thus $\Psi=g^*$. The map $g$ is unique because its component functions are
  the images of the restricted coordinates.

  Conversely, pullback along a smooth map $g\colon Y\to X$ preserves every
  smooth operation pointwise, and hence defines a $\Cinfty$-homomorphism. The
  two constructions are inverse.
\end{proof}

The corresponding result for all locally closed subsets is
Moerdijk--Reyes, Proposition~1.5
\cite[Chapter~I, Proposition~1.5]{Moerdijk_Reyes_Smooth_Infinitesimal_Analysis}.

For the usual class of finite-dimensional, second-countable smooth manifolds,
Whitney's embedding theorem may be strengthened to give a proper Euclidean
embedding \cite{lee2000smooth}. Its image is closed. The intrinsic smooth
functions and smooth maps agree with those in
\Cref{def:Smooth_Functions_On_Subsets}. We therefore obtain the case which
will be used later.

\begin{corollary}[Recovering manifolds]
  \label[corollary]{cor:Recovering_Manifolds}
  The contravariant functor
  \[
    \Cinfty(-)\colon
    (\text{smooth manifolds})\catop
    \longrightarrow
    \Cinfty\text{-}\mathrm{Ring}
  \]
  is fully faithful.
\end{corollary}

The conclusion is stronger than a reconstruction of points. The entire smooth
map $Y\to X$ is recovered from the homomorphism on function rings by looking
at the images of finitely many ambient coordinate functions.

\section{Summary and supplementary materials}
\label[section]{sec:Cinfty_Summary_And_Further_Reading}

\subsection{The live route}
\label[subsection]{sec:Cinfty_Live_Route}

The default live route occupies 70 minutes. Its engine is the free
$\Cinfty$-ring theorem together with Hadamard's lemma and the quotient
construction. The dual numbers are the principal calculation. Full
faithfulness for closed subsets is the geometric payoff and the handoff to
Talk 3. The locally closed extension and Borel's theorem remain available as
supplementary reading, but no later talk depends on their proofs.

\subsection{Takeaways and next step}
\label[subsection]{sec:Cinfty_Takeaways_And_Next_Step}

The main route through the chapter can be compressed into four statements.

\begin{enumerate}
  \item A $\Cinfty$-ring is an algebraic object in which every smooth map
    $\R^n\to\R$ acts as an operation.
  \item The ring $\Cinfty(\R^n)$ is free on its coordinate functions, and
    Hadamard's lemma implies that every ordinary ideal defines a quotient
    $\Cinfty$-ring.
  \item The quotients by $(t)$ and $(t^2)$ have the same real point but
    different algebraic structures. The latter retains a first-order
    thickening which the zero set forgets.
  \item Passing contravariantly from a manifold to its $\Cinfty$-ring of
    smooth functions loses neither the manifold nor its smooth maps.
\end{enumerate}

We have therefore found a classical algebraic language for smooth equations.
Two important tasks remain before deriving it. First, we need to understand
which $\Cinfty$-rings actually arise from manifolds and how geometric products
and transverse pullbacks appear algebraically. Second, ordinary quotients do
not yet have the homotopy-invariant universal property required for
nontransverse intersections. The next chapter addresses the first task and
prepares the second.

\subsection{Supplement: infinite jets and Borel's theorem}
\label[subsection]{sec:How_Large_Smooth_Algebra_Can_Be}

Borel's theorem is a useful calibration of the size of smooth algebra, but no
later chapter depends on it. The preceding examples were finite-dimensional
as ordinary $\R$-algebras. The next theorem shows how misleading that special
feature can be. Let
$\mathfrak m_0^\infty\subseteq\Cinfty(\R^n)$ be the ideal of functions
\emph{flat at the origin}, meaning that every partial derivative of every
order vanishes at the origin. The Taylor homomorphism is
\[
  T_0\colon\Cinfty(\R^n)\longrightarrow\R[[x_1,\ldots,x_n]],
  \qquad
  f\longmapsto
  \sum_{\alpha\in\N^n}\frac{D^\alpha f(0)}{\alpha!}x^\alpha.
\]

\begin{theorem}[Borel's theorem]
  \label[theorem]{thm:Borel}
  The Taylor homomorphism $T_0$ is surjective and has kernel
  $\mathfrak m_0^\infty$. It therefore induces an isomorphism of commutative
  $\R$-algebras
  \[
    \Cinfty(\R^n)/\mathfrak m_0^\infty
    \cong
    \R[[x_1,\ldots,x_n]].
  \]
  In particular, the formal power series ring inherits a $\Cinfty$-ring
  structure from the quotient on the left.
\end{theorem}

This is the form of Borel's theorem recorded as Moerdijk--Reyes,
Theorem~1.3
\cite[Chapter~I, Theorem~1.3]{Moerdijk_Reyes_Smooth_Infinitesimal_Analysis}.

The assertion about the kernel is immediate from the definition. Surjectivity
is the content of the theorem: every formal Taylor series is the Taylor series
of some smooth function. We record the standard construction in one variable
as a sketch, both to indicate why the result is plausible and to keep the
logical status of the argument clear.

\begin{discussion}[Proof sketch in one variable]
  \label[discussion]{disc:Borel_Proof_Sketch}
  Choose a compactly supported smooth function $\chi$ which is identically one
  near the origin. Given prescribed derivatives $(a_k)_{k\geq0}$, consider
  \[
    \sum_{k\geq0}\frac{a_k}{k!}t^k\chi(\lambda_k t).
  \]
  The numbers $\lambda_k$ are chosen recursively and sufficiently large. For
  every $j<k$, the $C^j$-norm of the $k$th summand can then be made at most
  $2^{-k}$. It follows that the series and each of its derivatives converge
  uniformly on $\R$. Near the origin the cutoff in the $k$th summand is
  constant, so this summand contributes $a_k$ to the $k$th derivative at zero
  and contributes nothing to the other Taylor coefficients. The resulting
  smooth function has the prescribed Taylor series. The multivariable proof
  uses the same rapidly shrinking cutoff construction, organized by total
  degree.
\end{discussion}

The quotient in \Cref{thm:Borel} is generated as a $\Cinfty$-ring by the
images of $x_1,\ldots,x_n$, but its underlying commutative algebra contains
arbitrary formal power series. This is the promised warning about finiteness:
generation under all smooth operations is a much weaker condition than
generation of the underlying commutative algebra.

\subsection{Supplement: locally closed subsets}
\label[subsection]{sec:Supplement_Locally_Closed_Subsets}

The live route through the talk can stop after
\Cref{cor:Recovering_Manifolds}. For completeness, this subsection proves the
extension used by Moerdijk and Reyes: the same full-faithfulness statement
holds for locally closed subsets. The only additional task is to replace an
open subset by a closed hypersurface in one higher-dimensional Euclidean
space. The construction follows their Lemma~1.4 and Proposition~1.5
\cite[Chapter~I, Lemma~1.4 and Proposition~1.5]{Moerdijk_Reyes_Smooth_Infinitesimal_Analysis}.

\begin{lemma}[Open subsets are smooth cozero loci]
  \label[lemma]{lem:Open_Subset_Cozero_Locus}
  For every open subset $U\subseteq\R^n$, there is a smooth function
  $f\colon\R^n\to[0,\infty)$ such that
  \[
    U=\{x\in\R^n\mid f(x)\neq0\}.
  \]
\end{lemma}

\begin{proof}
  Choose a countable collection of open balls $B_k$ which covers $U$ and whose
  closures lie in $U$. For each $k$, choose a nonnegative smooth function
  $\rho_k$ supported in $U$ and strictly positive on $B_k$. Choose constants
  $\varepsilon_k>0$ so small that
  \[
    \varepsilon_k
    \sup_{x\in\R^n}|D^\alpha\rho_k(x)|
    \leq 2^{-k}
  \]
  for every multi-index $\alpha$ with $|\alpha|\leq k$. Then
  \[
    f=\sum_{k\geq1}\varepsilon_k\rho_k
  \]
  converges together with all its derivatives and hence defines a smooth
  function. It vanishes outside $U$. At every point of $U$, at least one
  $\rho_k$ is positive, so $f$ is positive there.
\end{proof}

\begin{proposition}[Closed realization]
  \label[proposition]{prop:Locally_Closed_Closed_Realization}
  Every locally closed subset $X\subseteq\R^n$ is diffeomorphic, in the sense
  of \Cref{def:Smooth_Functions_On_Subsets}, to a closed subset of
  $\R^{n+1}$.
\end{proposition}

\begin{proof}
  First let $U\subseteq\R^n$ be open, and choose $f$ as in
  \Cref{lem:Open_Subset_Cozero_Locus}. The graph
  \[
    H_f
    =\{(x,y)\in\R^n\times\R\mid y f(x)=1\}
  \]
  is closed, since it is the zero locus of the continuous function
  $(x,y)\mapsto y f(x)-1$. The map
  \[
    U\longrightarrow H_f,
    \qquad
    x\longmapsto\bigl(x,1/f(x)\bigr)
  \]
  is a diffeomorphism with inverse the first projection.

  Now write the locally closed subset as $X=F\cap U$, with $F$ closed in
  $\R^n$ and $U$ open. Under the displayed diffeomorphism, $X$ corresponds to
  \[
    H_f\cap(F\times\R),
  \]
  which is closed in $\R^{n+1}$.
\end{proof}

\begin{corollary}[Moerdijk--Reyes]
  \label[corollary]{cor:Full_Faithfulness_Locally_Closed}
  The contravariant functor from locally closed subsets of Euclidean spaces and
  smooth maps to finitely generated $\Cinfty$-rings is fully faithful.
\end{corollary}

\begin{proof}
  Replace the source and target by the closed realizations from
  \Cref{prop:Locally_Closed_Closed_Realization}, and apply
  \Cref{thm:Full_Faithfulness_Closed_Subsets}. Their function rings are
  finitely generated because each closed realization is a quotient of a free
  $\Cinfty$-ring.
\end{proof}

\subsection{Related literature}
\label[subsection]{sec:Cinfty_Related_Literature}

The primary source for this chapter is Moerdijk--Reyes, Chapter~I, Section~1
\cite[Chapter~I, Section~1]{Moerdijk_Reyes_Smooth_Infinitesimal_Analysis}.
It contains the free-ring theorem, quotient construction, Borel's theorem, the
cozero-locus lemma, and full faithfulness for locally closed subsets. Our proof
of the closed-subset case makes the ambient-coordinate and bump-function
argument more explicit than the source.

Joyce gives a concise modern treatment of $C^\infty$-rings, finite generation
and presentation, Weil algebras, and the passage to $C^\infty$-schemes
\cite[Sections~2--3]{Joyce_Introduction_Cinfty_Schemes}. The latter passage is
deliberately postponed to Talk 7, after the universal construction, the first
affine derived calculations, and their stable linearization. Spivak's
construction shows how ordinary
$C^\infty$-rings are replaced by their homotopical counterparts in one model
of derived smooth manifolds
\cite[Sections~5--7]{Spivak_Derived_Smooth_Manifolds}.
